\begin{table}[htbp]
\centering
\scriptsize
\caption{Post-hoc comparison of clinical, imaging-report and pathology variables across the primary clusters.}
\label{tab:ue-posthoc}
\resizebox{\linewidth}{!}{%
\begin{tabular}{llrrrrr}
\toprule
Variable & Cluster 1 | 2 | 3 & Effect & $q$ & $p_{\text{adj}}$ & $q_{\text{adj}}$ & $p_{\text{C1}}$ \\
\midrule
LVI & 24\% $|$ 15\% $|$ 19\% & 0.090 ($V$) & 0.730 & $<$0.001 & 0.003 & 0.381 \\
ADC value (radiologist) & 0.75 [0.67-0.85] $|$ 0.83 [0.73-0.90] $|$ 0.82 [0.73-0.90] & 0.045 ($\epsilon^2$) & 0.001 & 0.015 & 0.154 & 0.003 \\
Maximum diameter (cm) & 3.35 [2.40-4.80] $|$ 3.70 [2.13-4.75] $|$ 2.95 [2.15-4.35] & 0.000 ($\epsilon^2$) & 0.829 & 0.655 & 0.794 & 0.585 \\
Deep myometrial invasion & 46\% $|$ 48\% $|$ 41\% & 0.046 ($V$) & 0.879 & 0.587 & 0.771 & 0.627 \\
Age (years) & 59.50 [52.00-67.00] $|$ 57.00 [52.00-67.50] $|$ 58.50 [52.00-66.00] & 0.000 ($\epsilon^2$) & 0.904 & 0.744 & 0.822 & 0.941 \\
Postmenopausal & 54\% $|$ 47\% $|$ 58\% & 0.069 ($V$) & 0.781 & 0.236 & 0.550 & 0.064 \\
Tumour margin & 78\% $|$ 82\% $|$ 84\% & 0.063 ($V$) & 0.781 & 0.906 & 0.919 & 0.981 \\
DWI hyperintensity & 94\% $|$ 91\% $|$ 94\% & 0.043 ($V$) & 0.879 & 0.176 & 0.529 & 0.226 \\
Ascites & 29\% $|$ 28\% $|$ 20\% & 0.069 ($V$) & 0.781 & 0.478 & 0.765 & 0.582 \\
Lymph-node positive & 18\% $|$ 11\% $|$ 9\% & 0.102 ($V$) & 0.730 & 0.338 & 0.680 & 0.355 \\
Cystic/solid composition & 23/43/33 $|$ 23/49/28 $|$ 25/47/28 & 0.044 ($V$) & 0.904 & 0.529 & 0.765 & 0.113 \\
Enhancement ratio, early arterial & 0.73 [0.61-0.88] $|$ 0.73 [0.63-0.97] $|$ 0.71 [0.59-0.88] & 0.000 ($\epsilon^2$) & 0.781 & 0.356 & 0.680 & 0.640 \\
Enhancement ratio, late arterial & 0.77 [0.65-0.94] $|$ 0.83 [0.67-0.94] $|$ 0.77 [0.64-0.92] & 0.002 ($\epsilon^2$) & 0.730 & 0.136 & 0.525 & 0.090 \\
Enhancement ratio, parenchymal & 0.84 [0.71-1.09] $|$ 0.82 [0.70-1.02] $|$ 0.98 [0.72-1.18] & 0.004 ($\epsilon^2$) & 0.730 & 0.150 & 0.525 & 0.516 \\
Enhancement ratio, delayed & 0.90 [0.71-1.09] $|$ 0.93 [0.72-1.11] $|$ 0.88 [0.67-1.09] & 0.000 ($\epsilon^2$) & 0.879 & 0.680 & 0.794 & 0.910 \\
Relative T2 signal & 1.47 [1.15-1.76] $|$ 1.43 [1.17-1.68] $|$ 1.56 [1.20-1.80] & 0.000 ($\epsilon^2$) & 0.781 & 0.230 & 0.550 & 0.269 \\
T2WI signal grade & 50\% $|$ 51\% $|$ 48\% & 0.013 ($V$) & 0.975 & 0.919 & 0.919 & 0.909 \\
p53 abnormal & 49\% $|$ 48\% $|$ 50\% & 0.011 ($V$) & 0.975 & 0.547 & 0.765 & 0.924 \\
Ki-67 (\%) & 45.00 [35.00-55.00] $|$ 35.00 [30.00-50.00] $|$ 40.00 [38.75-55.00] & 0.008 ($\epsilon^2$) & 0.724 & 0.036 & 0.196 & 0.008 \\
ER (\%) & 45.00 [35.00-55.00] $|$ 45.00 [37.50-60.00] $|$ 50.00 [40.00-60.00] & 0.003 ($\epsilon^2$) & 0.730 & 0.037 & 0.196 & 0.322 \\
PR (\%) & 52.50 [40.00-65.00] $|$ 50.00 [40.00-65.00] $|$ 55.00 [40.00-70.00] & 0.000 ($\epsilon^2$) & 0.879 & 0.447 & 0.765 & 0.339 \\
\bottomrule
\end{tabular}}
\par\smallskip\parbox{0.95\linewidth}{\footnotesize Continuous: median [IQR], Kruskal--Wallis; binary: \% positive; composition: \% per category; $\chi^2$ or Monte-Carlo $\chi^2$. adj: batch-adjusted likelihood-ratio test; C1: within center 1. $q$: BH-FDR across variables. No variable met the retention rule ($q<0.05$ unadjusted and adjusted with at least a small effect).}
\end{table}
